% GENERATED FILE. Edit canonical lessons, then run npm run generate:books.
% canonical-source-hash: fnv1a64:c069574d75f2931e
% canonical-lessons: PA-C54-dada, PA-C54-dadi, PA-C54-nana, PA-C54-nani, PA-C54-pati

\chapter{Grandfather, Grandmother, Grandfather, Grandmother, Husband}
\label{ch:pa-grandfather-grandmother-grandfather-grandmother-husband}
\hlchaptermodality{54}

\begin{chapteropening}
\textbf{By the end of this chapter:} \emph{I can say a grandfather, a grandmother, a grandfather, a grandmother and a husband.}

Complete the last lesson of chapter 54: I can say a grandfather, a grandmother, a grandfather, a grandmother and a husband.
\end{chapteropening}

\section[\texorpdfstring{\pa{ਦਾਦਾ} (dādā)}{ਦਾਦਾ (dādā)}]{\pa{ਦਾਦਾ} (dādā) — a grandfather (father's father)}

\label{lesson:PA-C54-dada}



\begin{quote}
Before the new one: say the Punjabi for a son, then the Punjabi for a daughter.
\end{quote}

\subsection*{You'll want to know: \pa{ਦਾਦਾ}}
\textbf{\pa{ਦਾਦਾ}} — \emph{dādā} — \textquotedblleft{}a grandfather (father's father)\textquotedblright{}.

\begin{cousinweb}
\textbf{\pa{ਦਾਦਾ}} (\emph{dādā}) is a tadbhava, a word inherited and reshaped by sound change. It comes from \emph{*dādda}, \textquotedblleft{}father or other elderly relative\textquotedblright{}, a form CDIAL reconstructs from its descendants (CDIAL 6261). CDIAL lists Punjabi \emph{dāddā}, \textquotedblleft{}father's father\textquotedblright{}.
\end{cousinweb}

\begin{grammarlens}[title={What you've built}]
One more word for this chapter.
\end{grammarlens}

\subsection*{Your turn}
\begin{itemize}
\raggedright
  \item \emph{Say it:} \emph{dādā}
  \item \emph{Say it:} \emph{dādā}, once more
  \item \emph{Say it:} say \emph{dādā}
  \item \emph{Recall:} say the Punjabi for a son, then the Punjabi for a daughter, then say \emph{dādā} again
\end{itemize}

\begin{tcolorbox}[breakable,colback=teal!4,colframe=teal!35!black,title={Before you move on}]
What does \pa{ਦਾਦਾ} mean? (\textquotedblleft{}A grandfather (father's father)\textquotedblright{}.) Say it once more.
\end{tcolorbox}
\section[\texorpdfstring{\pa{ਦਾਦੀ} (dādī)}{ਦਾਦੀ (dādī)}]{\pa{ਦਾਦੀ} (dādī) — a grandmother (father's mother)}

\label{lesson:PA-C54-dadi}



\begin{quote}
Before the new one: say the Punjabi for a daughter, then the Punjabi for a grandfather.
\end{quote}

\subsection*{You'll want to know: \pa{ਦਾਦੀ}}
\textbf{\pa{ਦਾਦੀ}} — \emph{dādī} — \textquotedblleft{}a grandmother (father's mother)\textquotedblright{}.

\begin{cousinweb}
\textbf{\pa{ਦਾਦੀ}} (\emph{dādī}) is a tadbhava, a word inherited and reshaped by sound change. It comes from \emph{*dādda}, \textquotedblleft{}father or other elderly relative\textquotedblright{}, a form CDIAL reconstructs from its descendants (CDIAL 6261). It is the feminine of \emph{dādā}; CDIAL lists Punjabi \emph{dāddī} beside \emph{dāddā}.
\end{cousinweb}

\begin{grammarlens}[title={What you've built}]
One more word for this chapter.
\end{grammarlens}

\subsection*{Your turn}
\begin{itemize}
\raggedright
  \item \emph{Say it:} \emph{dādī}
  \item \emph{Say it:} \emph{dādī}, once more
  \item \emph{Say it:} say \emph{dādī}
  \item \emph{Recall:} say the Punjabi for a daughter, then the Punjabi for a grandfather, then say \emph{dādī} again
\end{itemize}

\begin{tcolorbox}[breakable,colback=teal!4,colframe=teal!35!black,title={Before you move on}]
What does \pa{ਦਾਦੀ} mean? (\textquotedblleft{}A grandmother (father's mother)\textquotedblright{}.) Say it once more.
\end{tcolorbox}
\section[\texorpdfstring{\pa{ਨਾਨਾ} (nānā)}{ਨਾਨਾ (nānā)}]{\pa{ਨਾਨਾ} (nānā) — a grandfather (mother's father)}

\label{lesson:PA-C54-nana}



\begin{quote}
Before the new one: say the Punjabi for a grandfather, then the Punjabi for a grandmother.
\end{quote}

\subsection*{You'll want to know: \pa{ਨਾਨਾ}}
\textbf{\pa{ਨਾਨਾ}} — \emph{nānā} — \textquotedblleft{}a grandfather (mother's father)\textquotedblright{}.

\begin{cousinweb}
\textbf{\pa{ਨਾਨਾ}} (\emph{nānā}) is a tadbhava, a word inherited and reshaped by sound change. It comes from \emph{*nānna}, \textquotedblleft{}a term of respect for an older relative\textquotedblright{}, a form CDIAL reconstructs from its descendants (CDIAL 7059). CDIAL lists Punjabi \emph{nānā}, \textquotedblleft{}mother's father\textquotedblright{}.
\end{cousinweb}

\begin{grammarlens}[title={What you've built}]
One more word for this chapter.
\end{grammarlens}

\subsection*{Your turn}
\begin{itemize}
\raggedright
  \item \emph{Say it:} \emph{nānā}
  \item \emph{Say it:} \emph{nānā}, once more
  \item \emph{Say it:} say \emph{nānā}
  \item \emph{Recall:} say the Punjabi for a grandfather, then the Punjabi for a grandmother, then say \emph{nānā} again
\end{itemize}

\begin{tcolorbox}[breakable,colback=teal!4,colframe=teal!35!black,title={Before you move on}]
What does \pa{ਨਾਨਾ} mean? (\textquotedblleft{}A grandfather (mother's father)\textquotedblright{}.) Say it once more.
\end{tcolorbox}
\section[\texorpdfstring{\pa{ਨਾਨੀ} (nānī)}{ਨਾਨੀ (nānī)}]{\pa{ਨਾਨੀ} (nānī) — a grandmother (mother's mother)}

\label{lesson:PA-C54-nani}



\begin{quote}
Before the new one: say the Punjabi for a grandmother, then the Punjabi for a grandfather.
\end{quote}

\subsection*{You'll want to know: \pa{ਨਾਨੀ}}
\textbf{\pa{ਨਾਨੀ}} — \emph{nānī} — \textquotedblleft{}a grandmother (mother's mother)\textquotedblright{}.

\begin{cousinweb}
\textbf{\pa{ਨਾਨੀ}} (\emph{nānī}) is a tadbhava, a word inherited and reshaped by sound change. It comes from \emph{*nānna}, \textquotedblleft{}a term of respect for an older relative\textquotedblright{}, a form CDIAL reconstructs from its descendants (CDIAL 7059). It is the feminine of \emph{nānā}, \textquotedblleft{}mother's father\textquotedblright{}.
\end{cousinweb}

\begin{grammarlens}[title={What you've built}]
One more word for this chapter.
\end{grammarlens}

\subsection*{Your turn}
\begin{itemize}
\raggedright
  \item \emph{Say it:} \emph{nānī}
  \item \emph{Say it:} \emph{nānī}, once more
  \item \emph{Say it:} say \emph{nānī}
  \item \emph{Recall:} say the Punjabi for a grandmother, then the Punjabi for a grandfather, then say \emph{nānī} again
\end{itemize}

\begin{tcolorbox}[breakable,colback=teal!4,colframe=teal!35!black,title={Before you move on}]
What does \pa{ਨਾਨੀ} mean? (\textquotedblleft{}A grandmother (mother's mother)\textquotedblright{}.) Say it once more.
\end{tcolorbox}
\section[\texorpdfstring{\pa{ਪਤੀ} (patī)}{ਪਤੀ (patī)}]{\pa{ਪਤੀ} (patī) — a husband}

\label{lesson:PA-C54-pati}



\begin{quote}
Before the new one: say the Punjabi for a grandfather, then the Punjabi for a grandmother.
\end{quote}

\subsection*{You'll want to know: \pa{ਪਤੀ}}
\textbf{\pa{ਪਤੀ}} — \emph{patī} — \textquotedblleft{}a husband\textquotedblright{}.

\begin{cousinweb}
\textbf{\pa{ਪਤੀ}} (\emph{patī}) is Sanskrit \emph{pati}, \textquotedblleft{}master, husband\textquotedblright{}, taken over whole as a learned loan, a tatsama (CDIAL 7727). CDIAL lists the inherited Punjabi form as \emph{paī}, \textquotedblleft{}husband\textquotedblright{}.
\end{cousinweb}

\begin{grammarlens}[title={What you've built}]
One more word for this chapter.
\end{grammarlens}

\subsection*{Your turn}
\begin{itemize}
\raggedright
  \item \emph{Say it:} \emph{patī}
  \item \emph{Say it:} \emph{patī}, once more
  \item \emph{Say it:} say \emph{patī}
  \item \emph{Recall:} say the Punjabi for a grandfather, then the Punjabi for a grandmother, then say \emph{patī} again
\end{itemize}

\begin{tcolorbox}[breakable,colback=teal!4,colframe=teal!35!black,title={Before you move on}]
What does \pa{ਪਤੀ} mean? (\textquotedblleft{}A husband\textquotedblright{}.) Say it once more.
\end{tcolorbox}
